% THIS DOCUMENT IS FOLLOWS THE VOLERE TEMPLATE BY Suzanne Robertson and James Robertson
% ONLY THE SECTION HEADINGS ARE PROVIDED
%
% Initial draft from https://github.com/Dieblich/volere
%
% Risks are removed because they are covered by the Hazard Analysis
\documentclass[12pt]{article}

\usepackage{booktabs}
\usepackage{tabularx}
\usepackage{hyperref}
\hypersetup{
    bookmarks=true,         % show bookmarks bar?
      colorlinks=true,      % false: boxed links; true: colored links
    linkcolor=red,          % color of internal links (change box color with linkbordercolor)
    citecolor=green,        % color of links to bibliography
    filecolor=magenta,      % color of file links
    urlcolor=cyan           % color of external links
}

\newcommand{\lips}{\textit{Insert your content here.}}

\input{../Comments.text}
\input{../Common.text}

\begin{document}

\title{Software Requirements Specification for \progname: subtitle describing software} 
\author{\authname}
\date{\today}
	
\maketitle

~\newpage

\pagenumbering{roman}

\tableofcontents

~\newpage

\section*{Revision History}

\begin{tabularx}{\textwidth}{p{3cm}p{2cm}X}
\toprule {\textbf{Date}} & {\textbf{Version}} & {\textbf{Notes}}\\
\midrule
Date 1 & 1.0 & Notes\\
Date 2 & 1.1 & Notes\\
\bottomrule
\end{tabularx}

~\\

~\newpage
\section{Purpose of the Project}
\subsection{User Business}
\lips
\subsection{Goals of the Project}
\lips
\section{Stakeholders}
\subsection{Client}
\lips
\subsection{Customer}
\lips
\subsection{Other Stakeholders}
\lips
\subsection{Hands-On Users of the Project}
\lips
\subsection{Personas}
\subsubsection{Brandon - The Novice Saver and Avid Runner}

\subsubsection*{Background}
Brandon is a university student in his final year of Environmental Sciences at the University of Toronto, preparing for graduation in the spring. Brandon has recently begun learning more about his own finances, but still has a limited understanding of investing and saving.

\subsubsection*{Goals and Motivations}
Brandon enjoys recreational running as a way to stay active, relieve stress, and maintain a healthy lifestyle. He wants to improve running performance, build consistency, and take this hobby to the next level. At the same time, Brandon wants to develop his financial literacy and gain confidence in investment decisions before investing real savings.

\subsubsection*{Behaviours and Habits}
Brandon balances school responsibilities with running and personal development. He currently uses his phone as his main device to track his running activities such as goals and progress. Brandon prefers simple explanations, interactive learning, and measurable progress rather than complicated financial terminology.

\subsubsection*{Needs and Expectations}
\begin{itemize}
    \item A risk-free stock trading simulator that uses virtual money.
    \item Beginner-friendly lessons explaining investment opportunities, market trends, risk, and investment strategies.
    \item Progress tracking and feedback to help build investing knowledge and confidence.
    \item Running features to record workouts, monitor distance and pace, and track improvements over time.
    \item Achievable goals, milestones, and challenges that encourage consistent participation in both activities.
\end{itemize}

\subsubsection*{Frustrations}
Brandon finds financial platforms intimidating because of unfamiliar terminology and the potential to lose money. Many investing resources assume prior knowledge, while running apps may focus primarily on recording workouts rather than helping beginners improve. Brandon wants a straightforward experience that makes learning and progress rewarding.

\subsubsection{Connor – The Young Expert}
\subsubsection*{Background}
Connor is a third-year Commerce student specializing in Finance, returning to university after completing a year-long co-op placement at IG Wealth Management. Coming from a long line of financial professionals, Connor has developed a strong interest in financial management and aspires to build a successful career in the industry. As a top-performing student and hardworking individual, Connor is confident in his financial knowledge and enjoys applying what he has learned to real-world situations.

\subsubsection*{Goals and Motivations}
Connor wants to excel academically and build on the professional experience gained during his co-op placement to advance his career in financial management. He enjoys demonstrating his financial knowledge and sharing his opinions on investment opportunities and market trends. Although Connor is an avid sports fan, he recognizes that balancing academics, career ambitions, and personal interests leaves him with little time to exercise. He wants to find ways to incorporate physical activity into his busy schedule without compromising his professional goals.

\subsubsection*{Behaviours and Habits}
Connor maintains a demanding schedule focused on academics, career development, and expanding his financial knowledge. He regularly follows financial news, monitors market trends, and enjoys discussing investment strategies with others. Connor is competitive and motivated by measurable results, academic achievement, and opportunities to demonstrate his expertise. While he enjoys watching sports, exercise is not currently a consistent part of his routine due to his time commitments.

\subsubsection*{Needs and Expectations}
\begin{itemize}
    \item Advanced financial tools and simulated trading features that allow him to apply his existing investment knowledge.
    \item Access to market trends, investment opportunities, and performance statistics to support informed decision-making.
    \item Competitive features such as rankings, challenges, and performance comparisons to showcase his financial expertise.
    \item Efficient running and exercise tracking that fits into his busy academic and professional schedule.
    \item Personalized fitness goals, progress tracking, and achievable challenges that encourage consistent physical activity.
    \item An engaging platform that combines financial competition with physical activity, allowing him to pursue both interests in one place.
\end{itemize}

\subsubsection*{Frustrations}
Connor finds it difficult to balance his academic workload, career ambitions, and personal interests while maintaining an active lifestyle. Beginner-focused financial platforms may not provide enough depth or challenge to keep him engaged, while traditional fitness applications may not offer the competitive elements that motivate him. Connor wants an efficient, challenging experience that allows him to demonstrate his financial knowledge while encouraging him to make time for exercise.
\subsubsection{Daniella — The All-Arounder}
\subsubsection*{Background}
Daniella is a recent Civil Engineering graduate from the University of Guelph who gained valuable work experience and earned a substantial amount of money through multiple co-op placements throughout his degree. Through personal curiosity and working with her bank to understand her finances, Daniella developed a strong awareness of saving and managing money. Following graduation, she dedicated significant time to building her own investment portfolio and developing her financial knowledge. Although she has gained practical experience managing her savings, Daniella wants to become more confident in making independent investment decisions.

\subsubsection*{Goals and Motivations}
Daniella enjoys staying active through a variety of organized sports which he participates in throughout the year. She worries that the time commitment required by multiple sports leagues may eventually become unsustainable, and makes the decision to cut back on some of the organized leagues that she participates in. She is interested in beginning a more traditional circuit-style workout as these tend to be more convenient. Financially, she wants to continue growing her investment portfolio while gaining more experience evaluating opportunities and making important investment decisions independently.

\subsubsection*{Behaviours and Habits}
Daniella balances her professional responsibilities and personal interests while maintaining an active social life through organized sports. She prefers physical activities that are engaging, competitive, and social. Financially, Daniella takes an active role in monitoring his savings and managing his investment portfolio. She is curious, self-motivated, and willing to learn independently, but recognizes that managing a portfolio with real money offers limited opportunities to experiment with unfamiliar strategies without risking her savings.

\subsubsection*{Needs and Expectations}
\begin{itemize}
    \item A realistic stock trading simulator that allows her to practice making independent investment decisions without risking real money.
    \item Investment tools that help evaluate market trends, compare opportunities, and understand the potential risks and outcomes of different strategies.
    \item Performance tracking and feedback to evaluate investment decisions and identify areas for improvement.
    \item Convenient fitness tracking and personalized activity goals that fit her professional responsibilities and sports league commitments.
    \item Flexible exercise recommendations that provide alternatives to traditional gym workouts and help maintain fitness during busy periods.
    \item An integrated platform that encourages both financial development and physical activity without requiring excessive time or effort.
\end{itemize}

\subsubsection*{Frustrations}
Daniella finds it challenging to balance his professional life, multiple sports leagues, and financial goals without over-committing her time. Although she enjoys sports, she recognizes that relying entirely on organized leagues may make it difficult to maintain a consistent fitness routine year-round. Similarly, while she has successfully developed her own investment portfolio, she wants more opportunities to test strategies and learn from her decisions without putting his savings at risk. Daniella wants a flexible and practical solution that supports her financial growth and physical well-being while fitting naturally into his lifestyle.

\lips
\subsection{Priorities Assigned to Users}
\lips
\subsection{User Participation}
\lips
\subsection{Maintenance Users and Service Technicians}
\lips

\section{Mandated Constraints}
\subsection{Solution Constraints}
\lips
\subsection{Implementation Environment of the Current System}
\lips
\subsection{Partner or Collaborative Applications}
\lips
\subsection{Off-the-Shelf Software}
\lips
\subsection{Anticipated Workplace Environment}
\lips
\subsection{Schedule Constraints}
\lips
\subsection{Budget Constraints}
\lips
\subsection{Enterprise Constraints}
\lips

\section{Naming Conventions and Terminology}
\subsection{Glossary of All Terms, Including Acronyms, Used by Stakeholders
involved in the Project}
\lips

\section{Relevant Facts And Assumptions}
\subsection{Relevant Facts}
\lips
\subsection{Business Rules}
\lips
\subsection{Assumptions}
\lips

\section{The Scope of the Work}
\subsection{The Current Situation}
\lips
\subsection{The Context of the Work}
\lips
\subsection{Work Partitioning}
\lips
\subsection{Specifying a Business Use Case (BUC)}
\lips

\section{Business Data Model and Data Dictionary}
\subsection{Business Data Model}
\lips
\subsection{Data Dictionary}
\lips

\section{The Scope of the Product}
\subsection{Product Boundary}
\lips
\subsection{Product Use Case Table}
\lips
\subsection{Individual Product Use Cases (PUC's)}
\lips

\section{Functional Requirements}
\subsection{Functional Requirements}
\lips

\section{Look and Feel Requirements}
\subsection{Appearance Requirements}
\lips
\subsection{Style Requirements}
\lips

\section{Usability and Humanity Requirements}
\subsection{Ease of Use Requirements}
\lips
\subsection{Personalization and Internationalization Requirements}
\lips
\subsection{Learning Requirements}
\lips
\subsection{Understandability and Politeness Requirements}
\lips
\subsection{Accessibility Requirements}
\lips

\section{Performance Requirements}
\subsection{Speed and Latency Requirements}
\lips
\subsection{Safety-Critical Requirements}
\lips
\subsection{Precision or Accuracy Requirements}
\lips
\subsection{Robustness or Fault-Tolerance Requirements}
\lips
\subsection{Capacity Requirements}
\lips
\subsection{Scalability or Extensibility Requirements}
\lips
\subsection{Longevity Requirements}
\lips

\section{Operational and Environmental Requirements}
\subsection{Expected Physical Environment}
\lips
\subsection{Wider Environment Requirements}
\lips
\subsection{Requirements for Interfacing with Adjacent Systems}
\lips
\subsection{Productization Requirements}
\lips
\subsection{Release Requirements}
\lips

\section{Maintainability and Support Requirements}
\subsection{Maintenance Requirements}
\lips
\subsection{Supportability Requirements}
\lips
\subsection{Adaptability Requirements}
\lips

\section{Security Requirements}
\subsection{Access Requirements}
\lips
\subsection{Integrity Requirements}
\lips
\subsection{Privacy Requirements}
\lips
\subsection{Audit Requirements}
\lips
\subsection{Immunity Requirements}
\lips

\section{Cultural Requirements}
\subsection{Cultural Requirements}
\lips

\section{Compliance Requirements}
\subsection{Legal Requirements}
\lips
\subsection{Standards Compliance Requirements}
\lips

\section{Open Issues}
\lips

\section{Off-the-Shelf Solutions}
\subsection{Ready-Made Products}
\lips
\subsection{Reusable Components}
\lips
\subsection{Products That Can Be Copied}
\lips

\section{New Problems}
\subsection{Effects on the Current Environment}
\lips
\subsection{Effects on the Installed Systems}
\lips
\subsection{Potential User Problems}
\lips
\subsection{Limitations in the Anticipated Implementation Environment That May
Inhibit the New Product}
\lips
\subsection{Follow-Up Problems}
\lips

\section{Tasks}
\subsection{Project Planning}
\lips
\subsection{Planning of the Development Phases}
\lips

\section{Migration to the New Product}
\subsection{Requirements for Migration to the New Product}
\lips
\subsection{Data That Has to be Modified or Translated for the New System}
\lips

\section{Costs}
\lips
\section{User Documentation and Training}
\subsection{User Documentation Requirements}
\lips
\subsection{Training Requirements}
\lips

\section{Waiting Room}
\lips

\section{Ideas for Solution}
\lips

\newpage{}
\section{Appendix --- Generative AI Use Disclosure}

\wss{ChatGPT (version 5) was used to brainstorm ideas for the template for this disclose section.}

\subsection{AI Tools Used}

\wss{Give the name of the AI tool(s) used and the version.}

\subsection{How and Where Used}

\wss{Explain what the tool(s) were used for.  Examples include: brainstorming,
explaining a technical concept, reviewing the document, generating or modifying
text, generating or modifying diagrams, identifying errors or omissions or
inconsistencies.}

\wss{For each of the uses, identify which parts of the document were affected.}

\wss{You do not need to report every interaction, but you should disclose the 
uses that materially contributed to your work}

\subsection{Verification of AI-Generated Content}

\wss{\begin{itemize}
    \item Describe how the team checked AI-generated or AI-assisted material
    before including it in the document.
    \item In particular, verify factual claims, technical assertions,
    requirements, calculations, references, and other information for which an
    AI system may produce incorrect or fabricated results.
    \item \textbf{The team is responsible for the accuracy, completeness,
consistency, and appropriateness of the submitted document, regardless of
whether material was generated by a human or an AI system.}
\end{itemize}}

\newpage{}

\section{Appendix --- Reflection}

\input{../Reflection.text}

\input{../SRS_Reflection.text}

\end{document}